\chapter{Modelagem e Implementação do Data Warehouse}
\label{cap:dw}

Este capítulo documenta as decisões técnicas na construção do DW proposto, correspondendo à fase de desenvolvimento do ciclo DSR adotado (Capítulo~\ref{cap:metodologia}). A exposição segue o fluxo dos dados: fontes públicas, extração, \textit{staging}, modelagem dimensional e camada analítica que alimenta os \textit{dashboards}.
\section{Visão Geral da Arquitetura}

O sistema se organiza em camadas, mostradas na Figura~\ref{fig:arquitetura-dw}: extração, \textit{staging}, \textit{warehouse} dimensional e consumo via \textit{dashboards}. Separar fisicamente essas camadas atende a uma exigência operacional concreta: cada uma tem função definida e pode ser reexecutada sem arrastar as demais. Coleta, transformações e apresentação ficam desacopladas. Portais municipais podem mudar \textit{layout}; quando isso ocorre, o ajuste tende a ficar confinado à extração e ao \textit{staging}, sem reescrever consultas analíticas que já alimentam os \textit{dashboards}.

% \begin{figure}[ht]
% \begin{center}
% \begin{tabular}{|l|}
% \hline
% \textbf{Fontes} \\
% Portal Transparência PJF (XLS/XLSX) $\cdot$ STN (Ementário XLSX) \\
% $\cdot$ LOA PJF (\texttt{Funcionais.zip}, PDF/OCR) \\
% \hline
% $\downarrow$ \\
% \hline
% \textbf{Extração e Orquestração} (Dagster + Python) \\
% \texttt{despesa\_mensal\_consolidada}, \texttt{receita\_mensal\_comparativa}, \\
% \texttt{receita\_mensal\_prevista}, \texttt{stn\_ementario\_natureza\_receita}, \\
% \texttt{pjf\_loa\_funcao}, \texttt{pjf\_loa\_subfuncao} \\
% \hline
% $\downarrow$ \\
% \hline
% \textbf{Staging} (DuckDB, \textit{schema} \texttt{stg}, dbt views) \\
% Deduplicação, normalização leve, separação por layout \\
% \hline
% $\downarrow$ \\
% \hline
% \textbf{Warehouse Dimensional} (DuckDB, \textit{schema} \texttt{dwh}, dbt tables) \\
% 10 dimensões + 3 tabelas de fatos (modelagem Kimball) \\
% \hline
% $\downarrow$ \\
% \hline
% \textbf{Consumo} (Streamlit + Plotly) \\
% 7 páginas (5 analíticas + Início + Glossário), filtros globais \\
% \hline
% \end{tabular}
% \caption{Camadas do Data Warehouse construído para análise da execução orçamentária de Juiz de Fora.}
% \label{fig:arquitetura-dw}
% \end{center}
% \end{figure}


\begin{figure}[ht]
\centering
\begin{tikzpicture}[
    >={Stealth[length=3mm, width=2.5mm]},
    layer/.style={
        rectangle,
        rounded corners=3pt,
        draw=black!70,
        thick,
        minimum width=12cm,
        minimum height=1.35cm,
        align=center,
        inner sep=8pt
    },
    arrow/.style={
        ->,
        very thick,
        shorten >=3pt,
        shorten <=3pt
    },
    schemalabel/.style={
        font=\footnotesize\itshape,
        anchor=west,
        text=black!60
    }
]

% Camada 1: Fontes Públicas
\node[layer, fill=gray!8] (fontes) {%
    \textbf{Fontes Públicas}\\[2pt]
    \footnotesize Portal da Transparência PJF (XLS/XLSX) $\cdot$ STN Ementário (XLSX) $\cdot$ LOA PJF (PDF/OCR)%
};

% Camada 2: Extração
\node[layer, fill=gray!15, below=0.5cm of fontes] (extracao) {%
    \textbf{Extração e Orquestração}\\[2pt]
    \footnotesize Dagster $+$ Python $+$ \texttt{pandas} \quad $\longrightarrow$ \quad 6 \textit{assets} particionados%
};

% Camada 3: Staging
\node[layer, fill=gray!22, below=0.5cm of extracao] (staging) {%
    \textbf{Camada de \textit{Staging}} \quad \textnormal{\small(\textit{schema} \texttt{stg})}\\[2pt]
    \footnotesize DuckDB $+$ dbt \textit{views} \quad $\longrightarrow$ \quad deduplicação, normalização, separação por \textit{layout}%
};

% Camada 4: Warehouse Dimensional
\node[layer, fill=gray!30, below=0.5cm of staging] (warehouse) {%
    \textbf{\textit{Warehouse} Dimensional} \quad \textnormal{\small(\textit{schema} \texttt{dwh})}\\[2pt]
    \footnotesize DuckDB $+$ dbt \textit{tables} \quad $\longrightarrow$ \quad 10 dimensões $+$ 3 tabelas de fatos (modelagem Kimball)%
};

% Camada 5: Consumo
\node[layer, fill=gray!38, below=0.5cm of warehouse] (consumo) {%
    \textbf{Consumo Analítico}\\[2pt]
    \footnotesize Streamlit $+$ Plotly \quad $\longrightarrow$ \quad 7 páginas (5 analíticas $+$ Início $+$ Glossário)%
};

% Setas de fluxo
\draw[arrow] (fontes) -- (extracao);
\draw[arrow] (extracao) -- (staging);
\draw[arrow] (staging) -- (warehouse);
\draw[arrow] (warehouse) -- (consumo);

\end{tikzpicture}
\caption{Arquitetura em camadas do \textit{Data Warehouse} construído para análise da execução orçamentária de Juiz de Fora. As setas indicam o fluxo dos dados, das fontes públicas até o consumo analítico, seguindo o padrão \textit{medallion} adaptado.}
\label{fig:arquitetura-dw}
\end{figure}

A arquitetura segue de perto o padrão \textit{medallion} (\textit{bronze}/\textit{silver}/\textit{gold}) com adaptações. Os arquivos brutos ficam em \texttt{data/}, fora do banco. Tabelas brutas e \textit{views} de \textit{staging} compartilham o \textit{schema} \texttt{stg}: Dagster escreve as tabelas e dbt gera as views. O \textit{warehouse} dimensional ocupa o \textit{schema} \texttt{dwh}. A divisão por \textit{schema} torna explícita a fronteira entre dado bruto e dado modelado. Quem inspeciona o banco sabe imediatamente o que veio do portal e o que passou por transformações. Alterar uma modelo do dbt não obriga a baixar de novo planilhas do portal.

Todo o código que sustenta as decisões discutidas neste capítulo está versionado em um repositório público (\url{https://github.com/antoniomarcossouza/UFJF-DCC-TCC}), conforme o compromisso de reprodutibilidade firmado no Capítulo~\ref{cap:metodologia}.

\section{Tecnologias Utilizadas}

A \textit{stack} foi montada para rodar em máquina local e admitir \textit{deploy} simples, reprodutibilidade local e publicação no Streamlit Community Cloud, com ferramentas já difundidas em engenharia de dados. Para um trabalho com tempo limitado e uma pessoa na implementação, reprodutibilidade foi considerada uma característica importante: quem for replicar o trabalho precisa subir o ambiente sem contratar infraestrutura dedicada. A Tabela~\ref{tab:stack} resume os componentes.

\begin{table}[htb]
\caption{\textit{Stack} técnica utilizada no DW}
\label{tab:stack}
\centering
\small
\begin{tabular}{p{3.0cm} p{2.0cm} p{8.0cm}}
\hline
\textbf{Componente} & \textbf{Versão} & \textbf{Função no projeto} \\
\hline
Python & 3.13+ & Linguagem de extração e código geral \\
Dagster & 1.12 & Orquestrador de \textit{assets} \\
\texttt{pandas} & 3.0 & Leitura e normalização de planilhas \\
DuckDB & 1.5 & Banco analítico embarcado, formato OLAP colunar \\
dbt-core & 1.11 & Modelagem dimensional em SQL \\
dbt-duckdb & 1.10 & Adaptador dbt para DuckDB \\
\texttt{dbt\_utils} & 1.3 & Macros para chaves substitutas e \textit{date spine} \\
Streamlit & 1.57 & Construção dos \textit{dashboards} \\
Plotly & 6.7 & Visualizações interativas \\
\hline
\end{tabular}
\end{table}

DuckDB se trata de um banco analítico em processo, sem servidor, com armazenamento colunar voltado a agregações. Para escopo municipal e volume na casa de milhões de linhas, a escolha se apoia em fatores práticos; não há servidor para instalar, o banco é o arquivo \texttt{data/execucao\_orcamentaria.duckdb}. Consultas analíticas competem, neste porte, com opções pagas. A ponte com \texttt{pandas} e Arrow encurta o caminho entre Python e SQL, o que importa quando a extração já roda em \texttt{pandas} e a modelagem, em dbt. Como limitação, destaca-se a ausência de suporte robusto à concorrência multiusuário em ambientes de produção. Enquanto o consumo permanecesse restrito a uso local ou a poucos analistas, essa característica tendia a ter impacto limitado; com a publicação no Streamlit Community Cloud, o volume simultâneo de leitores passa a depender dos limites da plataforma e do arquivo DuckDB empacotado na aplicação.

Dagster e dbt dividem o trabalho de forma comum na engenharia de dados atual: o primeiro orquestra \textit{assets} de forma declarativa, o segundo concentra transformações em SQL. A integração via \texttt{dagster-dbt} amarra as duas partes. O \textit{asset} \texttt{dbt\_non\_partitioned\_models} só roda depois dos \textit{assets} de extração, de modo que modelagem e carga chegam juntas a cada ciclo. Na prática, quando uma planilha quebra, reexecuta-se Dagster sem tocar no SQL do dbt; quando alguma transformação muda, alteram-se modelos sem reescrever o download.

\section{Extração: Portal da Transparência, STN e LOA PJF}
\label{sec:extracao}

O Portal da Transparência da Prefeitura de Juiz de Fora não expõe API REST. A coleta depende de \texttt{HTTP GET} sobre planilhas Excel em URLs previsíveis. Esse desenho é típico de portais municipais brasileiros: a transparência existe, mas no formato que a prefeitura já usa internamente (relatórios exportados do SIAFEM), não no formato que um pipeline automatizado preferiria. O município cumpre a Lei de Acesso à Informação; quem quer analisar em série histórica paga o custo de interpretar planilhas heterogêneas. Daí o padrão adotado: baixar arquivo, localizar cabeçalho, interpretar com \texttt{pandas}. A Tabela~\ref{tab:fontes} sintetiza as fontes consumidas.

\begin{table}[htb]
\caption{Fontes de dados consumidas pelo \textit{pipeline} de extração}
\label{tab:fontes}
\centering
\small
\begin{tabular*}{\textwidth}{@{\extracolsep{\fill}}p{2.3cm}cc p{8.8cm}@{}}
\hline
\textbf{\textit{Asset}} & \textbf{Origem} & \textbf{Partição} & \textbf{Conteúdo} \\
\hline
\begin{minipage}[t]{2.1cm}\raggedright\path{despesa_mensal_consolidada}\end{minipage} & PJF & Mensal & Empenhos, liquidações e pagamentos do mês \\
\begin{minipage}[t]{2.1cm}\raggedright\path{receita_mensal_comparativa}\end{minipage} & PJF & Mensal & Previsão e arrecadação acumuladas no mês \\
\begin{minipage}[t]{2.1cm}\raggedright\path{receita_mensal_prevista}\end{minipage} & PJF & Anual & Previsão mensal de receitas para o exercício \\
\begin{minipage}[t]{2.1cm}\raggedright\path{stn_ementario_natureza_receita}\end{minipage} & STN & Anual & Códigos e descrições oficiais de natureza de receita \\
\begin{minipage}[t]{2.1cm}\raggedright\path{pjf_loa_funcao}\end{minipage} & LOA PJF & Sem partição & Funções orçamentárias (DimLOA, \texttt{Funcao.pdf}, exercício 2026) \\
\begin{minipage}[t]{2.1cm}\raggedright\path{pjf_loa_subfuncao}\end{minipage} & LOA PJF & Sem partição & Subfunções orçamentárias (DimLOA, \texttt{SubFuncao.pdf}, exercício 2026) \\
\hline
\end{tabular*}
\end{table}

As seguintes decisões de extração mostram como o pipeline reage a problemas reais de fonte pública, não a um \textit{dataset} estável de laboratório. Cada uma nasceu de uma falha durante a implementação.

\subsection{Detecção Dinâmica de Cabeçalho}

As planilhas de despesa não fixam a linha de cabeçalho. O arquivo abre com metadados (logotipo, título, período) antes da tabela. A extração percorre as primeiras linhas até achar a célula \texttt{"Unidade Administrativa"} e só então define o início dos dados. Sem essa varredura, um mês com uma linha a mais de cabeçalho derruba a carga inteira. Pequenas mudanças de \textit{layout} entre meses deixam de quebrar o pipeline.

\subsection{Quebra de \textit{Layout} em 2025-05}

A planilha de receita comparativa mudou de formato no arquivo de maio de 2025 (partição \texttt{2505}). Em vez de descartar histórico ou forçar um parser único, o código mantém dois leitores: \texttt{read\_receita\_comparativa\_pre2505} e \texttt{read\_receita\_comparativa\_2505}. A escolha ocorre na extração, pela chave de partição; assim, anos anteriores continuam reprocessáveis sem perda.

\subsection{Classificação Funcional da LOA (exercício 2026)}

A LOA PJF publica a classificação funcional (DimLOA) em \texttt{Funcionais.zip}, com um arquivo PDF por tabela (\texttt{Funcao.pdf}, \texttt{SubFuncao.pdf}). Os PDFs são imagens escaneadas, não texto selecionável. A extração converte cada página com \texttt{pdftoppm} (Poppler) e reconhece o texto com Tesseract (pacotes instalados na imagem Docker). O \textit{parser} valida o layout do exercício 2026 (29 funções, 117 subfunções) e grava duas tabelas brutas em \texttt{stg}. Cada exercício novo pode exigir recalibrar o reconhecimento.

Os \textit{assets} não são particionados: o exercício 2026 está fixo na constante \texttt{LOA\_YEAR}. Cada ano orçamentário traz PDFs com \textit{layout} distinto; generalizar a ingestão exige validar e, se necessário, adaptar o \textit{parser} ano a ano. Por ora, função e subfunção alimentam \texttt{dim\_funcao} e \texttt{dim\_subfuncao}, usadas no \textit{dashboard} para rotular cortes por área de governo. Não há dotação autorizada nesta carga.

\subsection{Idempotência por Estratégia de Escrita}

Cada \textit{asset} usa a estratégia de escrita que combina com sua semântica. Despesa e receita comparativa entram por \textit{append}; a deduplicação fica no \textit{staging}, por \texttt{nm\_arquivo} e \texttt{dt\_atualizacao}. O ementário da STN usa \textit{overwrite by partition}, substituindo todas as linhas do ano a cada execução. As tabelas da LOA, anuais e canônicas, usam sobrescrita total da tabela (\textit{overwrite}), pois carregam um único exercício por materialização. Os caminhos diferem, mas o efeito desejado é o mesmo: reexecutar o pipeline devolve o mesmo estado final do \textit{warehouse}; um \textit{backfill} de março de 2023 não deixa duplicata nem lacuna no conjunto de dados.

\section{Camada de Staging}

Seis \textit{views} dbt cobrem as seis fontes brutas (despesa, receita comparativa, receita prevista, ementário STN, função LOA, subfunção LOA). Há deduplicação por \texttt{nm\_arquivo} e \texttt{dt\_atualizacao} nas planilhas PJF; remoção de \texttt{TOTAIS GERAIS} na receita comparativa; normalização de códigos de natureza de receita (só dígitos, sem sufixos como \texttt{.0}) e \textit{trim} nas tabelas LOA. Regras mais pesadas (por exemplo, \textit{unpivot} da previsão mensal) ficam em modelos \textit{intermediate} efêmeros, entre \textit{staging} e \textit{warehouse}. Manter o \textit{staging} enxuto evita que uma \textit{view} vire depósito de transformações difíceis de testar.

% A deduplicação de despesa mostra por que \textit{append} na extração exige cuidado; O mesmo arquivo pode entrar várias vezes, sobretudo em \textit{backfills}. A \textit{view} de \textit{staging} resolve com janela:

% \begin{singlespacing}
% \begin{lstlisting}[language=SQL,basicstyle=\ttfamily\small,frame=single,breaklines=true]
% select *
% from {{ source('stg', 'pjf_despesa_mensal_consolidada') }}
% qualify rank() over (
%     partition by nm_arquivo
%     order by dt_atualizacao desc
% ) = 1
% \end{lstlisting}
% \end{singlespacing}

% A cláusula \texttt{qualify}, suportada pelo DuckDB e por dialetos SQL recentes, filtra o resultado de funções de janela sem subconsulta. Para cada arquivo de origem, permanece só a versão mais recente carregada. A vantagem prática é legibilidade: quem revisa o SQL enxerga a regra de deduplicação num único bloco, sem aninhar \texttt{SELECT} dentro de \texttt{SELECT}.

\section{Modelo Dimensional}

A modelagem segue a metodologia de Kimball: processos de negócio, granularidade, dimensões e fatos. Dois processos foram tratados em separado: execução de despesa e execução de receita. Tentar unificá-los num único fato misturaria grãos, dimensões e calendários de publicação distintos. Despesas e receitas no portal PJF saem em arquivos diferentes, com colunas e periodicidade próprias; forçar um fato único criaria colunas esparsas e consultas analíticas mais frágeis.

\subsection{Tabelas de Fatos}

O modelo conta com três tabelas de fatos, descritas na Tabela~\ref{tab:fatos}.

\begin{table}[htb]
\caption{Tabelas de fatos do \textit{warehouse}}
\label{tab:fatos}
\centering
\small
\begin{tabular*}{\textwidth}{@{\extracolsep{\fill}}p{2.2cm}p{7.2cm}p{5.2cm}@{}}
\hline
\textbf{Tabela} & \textbf{Granularidade} & \textbf{Medidas} \\
\hline
\begin{minipage}[t]{3.0cm}\raggedright\path{fct_despesa}\end{minipage} & Linha do arquivo mensal consolidado (combinação de empenho, unidade, fornecedor, funcional, natureza, fonte e datas) & \path{vl_empenhado_mes}, \path{vl_liquidado_mes}, \path{vl_pago_mes} \\
\begin{minipage}[t]{3.0cm}\raggedright\path{fct_receita}\end{minipage} & Natureza de receita por mês de referência & \path{vl_previsto_mensal}, \path{vl_arrecadada_mes} \\
\begin{minipage}[t]{3.0cm}\raggedright\path{fct_receita_acumulada}\end{minipage} & Natureza de receita por mês de referência & \path{vl_previsao_inicial_comparativa}, \path{vl_previsao_atualizada}, \path{vl_arrecadada_ano}, \path{vl_a_realizar} \\
\hline
\end{tabular*}
\end{table}

Uma observação sobre a granularidade de \texttt{fct\_despesa} é que, no modelo ideal, cada linha seria um empenho ou um estágio (empenho, liquidação, pagamento). O portal, porém, publica relatórios mensais já consolidados: cada linha agrega valores por combinação de dimensões. Modelar nesse grão respeita o que a fonte entrega, sem simular atomicidade inexistente.

\texttt{fct\_receita} e \texttt{fct\_receita\_acumulada} também nascem de semânticas diferentes. O arquivo de previsão mensal traz valor previsto por mês, alimentando \texttt{fct\_receita}. O comparativo traz acumulados no ano, alimentando \texttt{fct\_receita\_acumulada}. Juntar tudo numa tabela única exigiria misturar medidas de ritmo mensal com medidas de saldo anual.

\subsection{Dimensões}

O modelo inclui dez dimensões, descritas na Tabela~\ref{tab:dimensoes}.

\begin{table}[htb]
\caption{Dimensões do \textit{warehouse}}
\label{tab:dimensoes}
\centering
\small
\begin{tabular}{p{5.5cm} p{8.5cm}}
\hline
\textbf{Dimensão} & \textbf{Conteúdo principal} \\
\hline
\texttt{dim\_tempo} & Dia, mês, trimestre, semestre, ano (\textit{date spine} de 2014 até 1 ano à frente) \\
\texttt{dim\_unidade\_administrativa} & Órgão executor da despesa \\
\texttt{dim\_funcional\_pragmatica} & Código e descrição da funcional programática \\
\texttt{dim\_natureza\_despesa} & Código e descrição da natureza de despesa (execução PJF) \\
\texttt{dim\_fonte\_recurso} & Código e descrição da fonte de recurso \\
\texttt{dim\_fornecedor} & CPF/CNPJ e nome do fornecedor \\
\texttt{dim\_empenho} & Número da nota, modalidade, licitação, referência legal, processo e descrição \\
\texttt{dim\_natureza\_receita} & Código e descrição da natureza de receita (enriquecida pelo ementário STN) \\
\texttt{dim\_funcao} & Código e descrição de função (LOA PJF, exercício 2026) \\
\texttt{dim\_subfuncao} & Código e descrição de subfunção (LOA PJF, exercício 2026) \\
\hline
\end{tabular}
\end{table}

\texttt{dim\_natureza\_receita} cruza duas origens. Quando o ementário da STN cobre o código, prevalece a descrição oficial; nos demais casos, entra a descrição publicada pela PJF. Quem consulta o \textit{dashboard} vê nomenclatura nacional padronizada onde ela existe, sem apagar códigos que só aparecem no município.

\texttt{dim\_funcao} e \texttt{dim\_subfuncao} não se ligam diretamente aos fatos: a execução mensal traz a funcional programática composta em \texttt{dim\_funcional\_pragmatica}. O \textit{dashboard} extrai função e subfunção dessa \textit{string} e faz \textit{join} com as dimensões LOA para exibir nomes oficiais do exercício.

\subsection{Chaves Substitutas}

As dimensões usam chaves substitutas geradas por \texttt{dbt\_utils.generate\_surrogate\_key} (MD5 sobre concatenação determinística dos campos naturais). Nas dimensões de despesa derivadas do consolidado mensal, a chave nasce no modelo \textit{intermediate} e é reaproveitada no \textit{warehouse}; \texttt{dim\_tempo} usa a mesma macro sobre a data do dia. Kimball recomenda esse padrão por dois motivos operacionais. A chave não muda se um rótulo de exibição for corrigido na fonte. A junção entre fatos volumosos e dimensões usa um único \texttt{varchar(32)}. O hash determinístico garante que reprocessar o pipeline recria as mesmas chaves, preservando idempotência. Sem chave substituta, uma correção de grafia em \texttt{dim\_fornecedor} quebraria o histórico de junção ou exigiria atualização em cascata nos fatos.

\subsection{Esquema Estrela}

As ligações entre fatos e dimensões formam três tabelas de fatos articuladas por dimensões conformadas, com \texttt{dim\_tempo} como dimensão central e \texttt{dim\_natureza\_receita} compartilhada entre as duas tabelas de receita. A Figura~\ref{fig:estrelas} ilustra essas conexões.

\begin{figure}[ht]
\begin{center}
\begin{tabular}{|l|}
\hline
\textbf{Estrela de Despesa} \\
\texttt{fct\_despesa} $\leftarrow$ dim\_tempo (empenho, liquidação) \\
\texttt{fct\_despesa} $\leftarrow$ dim\_unidade\_administrativa \\
\texttt{fct\_despesa} $\leftarrow$ dim\_fornecedor \\
\texttt{fct\_despesa} $\leftarrow$ dim\_funcional\_pragmatica \\
\texttt{fct\_despesa} $\leftarrow$ dim\_natureza\_despesa \\
\texttt{fct\_despesa} $\leftarrow$ dim\_fonte\_recurso \\
\texttt{fct\_despesa} $\leftarrow$ dim\_empenho \\
\hline
\textbf{Estrela de Receita} \\
\texttt{fct\_receita} $\leftarrow$ dim\_tempo, dim\_natureza\_receita \\
\hline
\textbf{Estrela de Receita Acumulada} \\
\texttt{fct\_receita\_acumulada} $\leftarrow$ dim\_tempo, dim\_natureza\_receita \\
\hline
\end{tabular}
\caption{Esquemas estrela do \textit{warehouse}, compartilhando \texttt{dim\_tempo}.}
\label{fig:estrelas}
\end{center}
\end{figure}

\begin{figure}[ht]
\centering
\resizebox{\textwidth}{!}{
\begin{tikzpicture}[
    >={Stealth[length=2.5mm, width=2mm]},
    fact/.style={
        rectangle,
        rounded corners=2pt,
        draw=black!80,
        thick,
        fill=gray!30,
        minimum width=2.4cm,
        minimum height=0.9cm,
        align=center,
        font=\small\ttfamily
    },
    dim/.style={
        rectangle,
        rounded corners=2pt,
        draw=black!60,
        fill=gray!8,
        minimum width=2.4cm,
        minimum height=0.7cm,
        align=center,
        font=\footnotesize\ttfamily
    },
    dimshared/.style={
        rectangle,
        rounded corners=2pt,
        draw=black!80,
        thick,
        fill=gray!18,
        minimum width=2.4cm,
        minimum height=0.7cm,
        align=center,
        font=\footnotesize\ttfamily
    },
    link/.style={
        -,
        thick,
        black!60
    }
]

% =============== ESTRELA DE DESPESA (esquerda) ===============

% Fato central da despesa
\node[fact] (fct_despesa) at (0, 0) {fct\_despesa};

% Dimensões ao redor da despesa (não-compartilhadas)
\node[dim] (dim_unidade)     at (-3.6, 2.2) {dim\_unidade\_\\administrativa};
\node[dim] (dim_fornecedor)  at (-3.6, 0.6) {dim\_fornecedor};
\node[dim] (dim_funcional)   at (-3.6, -1.0) {dim\_funcional\_\\pragmatica};
\node[dim] (dim_nat_desp)    at (0, -2.4) {dim\_natureza\_\\despesa};
\node[dim] (dim_fonte)       at (0, 2.4) {dim\_fonte\_recurso};
\node[dim] (dim_empenho)     at (3.6, -1.0) {dim\_empenho};

% Ligações da despesa
\draw[link] (fct_despesa) -- (dim_unidade);
\draw[link] (fct_despesa) -- (dim_fornecedor);
\draw[link] (fct_despesa) -- (dim_funcional);
\draw[link] (fct_despesa) -- (dim_nat_desp);
\draw[link] (fct_despesa) -- (dim_fonte);
\draw[link] (fct_despesa) -- (dim_empenho);

% =============== DIMENSÕES CONFORMADAS (centro) ===============

% dim_tempo compartilhada — posicionada à direita da fct_despesa
\node[dimshared] (dim_tempo) at (4.8, 2.2) {dim\_tempo};

% Ligação dim_tempo <-> fct_despesa, com rótulo dos dois papéis sobre a aresta
\draw[link] (fct_despesa) -- node[
    midway,
    sloped,
    above,
    font=\scriptsize\itshape,
    text=black!60
] {(empenho, liquidação)} (dim_tempo);

% =============== ESTRELAS DE RECEITA (direita) ===============

% Fatos de receita
\node[fact] (fct_receita)     at (8.6, 1.0) {fct\_receita};
\node[fact] (fct_rec_acum)    at (8.6, -1.5) {fct\_receita\_\\acumulada};

% Dimensão compartilhada entre as duas receitas
\node[dimshared] (dim_nat_rec) at (12.4, -0.25) {dim\_natureza\_\\receita};

% Ligações das receitas
\draw[link] (fct_receita) -- (dim_tempo);
\draw[link] (fct_rec_acum) -- (dim_tempo);
\draw[link] (fct_receita) -- (dim_nat_rec);
\draw[link] (fct_rec_acum) -- (dim_nat_rec);

% =============== LEGENDA ===============

\begin{scope}[on background layer]
    \node[
        draw=black!40,
        dashed,
        rounded corners=3pt,
        fill=white,
        inner sep=4pt,
        font=\scriptsize
    ] at (5, -3.6) {%
        \begin{tabular}{@{}l@{\;\;}l@{}}
            \tikz\node[fact, minimum width=0.5cm, minimum height=0.25cm, font=\tiny]{}; & Tabela de fatos \\
            \tikz\node[dim, minimum width=0.5cm, minimum height=0.25cm, font=\tiny]{}; & Dimensão exclusiva \\
            \tikz\node[dimshared, minimum width=0.5cm, minimum height=0.25cm, font=\tiny]{}; & Dimensão conformada \\
        \end{tabular}
    };
\end{scope}

\end{tikzpicture}
}
\caption{Modelo dimensional do \textit{warehouse}: três tabelas de fatos articuladas por dimensões conformadas. \texttt{dim\_tempo} é compartilhada pelas três fatos e atua em dois papéis (\textit{role-playing}) na estrela de despesa. \texttt{dim\_natureza\_receita} é conformada entre \texttt{fct\_receita} e \texttt{fct\_receita\_acumulada}.}
\label{fig:estrelas}
\end{figure}

A dimensão \texttt{dim\_tempo} aparece na estrela de despesa em dois papéis distintos: data de empenho e data de liquidação. Esse padrão, conhecido como \textit{role-playing dimension} \cite{kimball2013data}, evita duplicar a estrutura da dimensão temporal para cada relacionamento. Na prática, a mesma tabela física é referenciada por chaves substitutas distintas nos fatos, permitindo análises que filtram simultaneamente por mês de empenho e por mês de liquidação, sem comprometer a integridade do modelo.

Nenhuma dimensão implementa SCD (\textit{slowly changing dimension}). Todas são tipo 1: guardam só o estado mais recente da fonte. Para descrições de natureza no padrão brasileiro, isso costuma bastar, porque a STN revisa ementários em ciclos anuais e as mudanças são pontuais. Reorganizações administrativas do município, apesar de raras, ficariam mal registradas: se uma secretaria mudar de nome, o histórico apareceria com o rótulo novo. Capturar essas mudanças exigiria SCD tipo 2.

\section{Processos de ETL e Idempotência}

A unidade de idempotência é a partição temporal nas fontes particionadas: mês de competência nas planilhas PJF e ano de referência na STN. A LOA funcional não usa partição; a idempotência vem da sobrescrita total das tabelas brutas a cada materialização. Reexecutar o pipeline deve devolver o mesmo \textit{warehouse}, independentemente de quantas cargas anteriores existam. Três mecanismos sustentam essa propriedade. A estratégia de escrita por fonte (Seção~\ref{sec:extracao}) trata STN com sobrescrita por partição, LOA com sobrescrita total e planilhas PJF com \textit{append} mais deduplicação. As chaves substitutas determinísticas evitam linhas falsas em dimensões quando uma partição é recarregada: o mesmo conjunto de campos naturais gera o mesmo MD5. A materialização \textit{table}, e não incremental, nas tabelas de fatos faz o dbt reconstruir o conjunto inteiro a partir do \textit{staging} a cada execução. Para o volume municipal atual, o \textit{full refresh} não compromete desempenho e garante que a tabela final reflita exatamente o que veio do \textit{staging}.

\section{Qualidade de Dados e Testes}

Arquivos \texttt{\_schema.yml} especificam verificações de unicidade e de nulidade nas chaves das dimensões, de unicidade nas combinações de chaves substitutas nos fatos, de presença das medidas e de integridade referencial. Testes singulares em SQL complementam esse conjunto com invariantes de receita: a arrecadação acumulada no exercício não pode ser inferior ao acumulado dos meses anteriores, e o valor a realizar deve coincidir com a previsão atualizada menos o arrecadado no ano.

Em dados fiscais publicados, falha silenciosa custa caro: um valor negativo por erro de \textit{parsing} ou uma quebra de monotonicidade entre meses pode virar conclusão errada sobre execução orçamentária. O conjunto de testes não cobre tudo, mas ataca as falhas mais prováveis: chave duplicada em dimensão, medida monetária inválida, carga parcial entre meses e \textit{drift} de \textit{schema} após mudança de \textit{layout} na fonte. A Tabela~\ref{tab:dq} resume os testes singulares.

\begin{table}[htb]
\caption{Testes singulares de qualidade de dados}
\label{tab:dq}
\centering
\small
\begin{tabular*}{\textwidth}{@{\extracolsep{\fill}}p{5.6cm}p{8.8cm}@{}}
\hline
\textbf{Teste} & \textbf{Invariante verificada} \\
\hline
\begin{minipage}[t]{6.4cm}\raggedright\path{dq_fct_receita_acumulada_nao_menor_que_mes}\end{minipage} & Valor acumulado no ano deve ser maior ou igual à arrecadação no mês corrente \\
\begin{minipage}[t]{6.4cm}\raggedright\path{dq_fct_receita_acumulada_monotonicidade}\end{minipage} & Valor acumulado no ano deve ser não decrescente mês a mês, por natureza \\
\begin{minipage}[t]{5.4cm}\raggedright\path{dq_fct_receita_vl_a_realizar}\end{minipage} & Valor a realizar deve coincidir com previsão atualizada menos arrecadado no ano \\
\hline
\end{tabular*}
\end{table}

\section{Limitações da Implementação}

Registrar o que o sistema não faz é tão importante quanto descrever o que faz. Três limitações marcam o recorte atual, delimitando perguntas respondíveis hoje e orientam quem quiser reproduzir ou estender o trabalho.

\subsection{Granularidade da Despesa}

A fonte publica relatórios mensais consolidados, não transações individuais. Análises de empenho a empenho exigiriam outras bases. O modelo define o grão pela linha do arquivo, mas os valores em \texttt{fct\_despesa} são movimentações mensais por combinação de dimensões, não atos administrativos isolados.

\subsection{Ausência de Orçamento Autorizado}

O \textit{warehouse} traz empenhado, liquidado e pago, mas não a dotação autorizada na LOA. A ingestão atual da LOA cobre apenas a classificação funcional, não os valores autorizados por unidade ou natureza. Sem autorizado, indicadores como percentual de execução sobre o previsto na lei ficam indisponíveis.

\subsection{Classificação MCASP Funcional Incompleta}

A dimensão \texttt{dim\_funcional\_pragmatica} guarda o código composto da funcional programática como \textit{string} única nos fatos. Para rotular função e subfunção no \textit{dashboard}, o sistema consulta \texttt{dim\_funcao} e \texttt{dim\_subfuncao}, alimentadas pela LOA PJF do exercício 2026, com \textit{fallback} textual quando o código da execução não consta no ementário carregado.

%\subsubsection{Rotulagem LOA e \textit{fallbacks} na execução funcional}

A revisão de qualidade do \textit{dashboard} quantificou a cobertura dos rótulos LOA na página Execução Orçamentária. Entre os 126 pares distintos função/subfunção presentes na execução, 18 (14{,}3\,\%) exibem \textit{fallback} do tipo \texttt{Função XX} ou \texttt{Subfunção XXX}: o código aparece na movimentação mensal, mas não foi resolvido contra \texttt{dim\_funcao} ou \texttt{dim\_subfuncao}. Os pares de maior volume executado recebem nomes corretos da LOA; os \textit{fallbacks} concentram-se em combinações de menor pagamento acumulado. Ou seja, o problema afeta sobretudo fatias pequenas da execução, mas ainda compromete a legibilidade de cortes que podem ser relevantes para auditoria.

As seguintes melhorias aparecem como prioridade para reduzir \textit{fallbacks} sem alterar a interface analítica. Revisão manual ou semi-automática dos PDFs após OCR, com tabela de correções versionada (CSV ou \textit{seed} dbt) para códigos sistematicamente omitidos. Enriquecimento cruzado com a tabela MCASP da STN ou com códigos já presentes em \texttt{dim\_funcional\_pragmatica} na execução, usando a LOA PJF como fonte primária e o ementário nacional como reserva para códigos ausentes no PDF municipal. A longo prazo, solicitar à prefeitura publicação dos dados em formato tabular eliminaria a dependência de OCR.

Essas limitações delimitam perguntas respondíveis hoje e apontam extensões concretas: dotação autorizada na LOA, grão mais fino de despesa, desagregação funcional completa, \textit{parser} LOA adaptado a exercícios anteriores e redução dos \textit{fallbacks} funcionais descritos acima.
